\documentclass[11pt]{article}

\usepackage[margin=1in]{geometry}
\usepackage{booktabs}
\usepackage{graphicx}
\usepackage{amsmath}

\setlength{\emergencystretch}{2em}

\title{Pipeline Methodology and Results Report}
\author{}
\date{2026-08-01}

\begin{document}
\maketitle

\section{Scope}

This report describes how the current LLM-assisted entity resolution (ER) pipeline
works end to end, explains each attribute-selection methodology it implements, and
compares results across those methodologies on three benchmark datasets:
Fodors--Zagat, DBLP--ACM, and Amazon--Walmart.

All numbers reported here come from run logs that match the pipeline's current
configuration: SentenceBERT embedding model \texttt{BAAI/bge-large-en-v1.5}, matcher
LLM \texttt{google/gemini-2.5-flash} (via OpenRouter), and the per-dataset blocking
and adaptive hyperparameters currently defined in \texttt{code/main.py}. A number of
older logs were excluded because they do not reflect the current pipeline: runs
embedded with a superseded model (\texttt{all-mpnet-base-v2}), runs matched with a
different LLM (\texttt{deepseek/deepseek-v4-flash} or \texttt{google/gemini-2.0-flash-001}),
pre-refactor logs that predate the current run-summary schema, one truncated
Fodors--Zagat run that only produced 1{,}105 of 2{,}665 expected result pairs, and two
stale aggregate CSVs (\texttt{logs/experiment\_summary.csv} and
\texttt{logs/step1\_metrics\_data.csv}) that were both generated under different
blocking hyperparameters than the pipeline currently uses. Excluding these keeps
every number in this report directly comparable.

One further exclusion is specific to Fodors--Zagat. During analysis its
\texttt{class} attribute was found to be the ground-truth cluster identifier
(Section~\ref{sec:leakage}), which makes any result computed with that attribute in
scope uninterpretable. All Fodors--Zagat numbers in this report therefore come from
a second generation of runs executed with \texttt{--exclude-attributes class}; the
earlier Fodors--Zagat runs are retained only in Section~\ref{sec:leakage}, where
they are used to quantify the effect of the leak. DBLP--ACM was checked and is
clean. Amazon--Walmart's \texttt{original\_id} attribute was also checked and does
\emph{not} leak (it disagrees on 100\% of ground-truth pairs), though it is
uninformative by construction and is a candidate for exclusion in future runs.

A third exclusion applies to DBLP--ACM. Of the thirteen runs in the reseeding sweep
attempted on 2 August, eight failed on every LLM call and a ninth failed on 28.9\% of
them, the API budget having been exhausted mid-sweep. Because the pipeline scores a
failed call as a non-match and still writes complete metrics, those nine runs report
numerically valid but meaningless results; all nine are excluded, leaving the five
\texttt{random} runs analysed in Section~\ref{sec:dblp-seeds}. A later Fodors--Zagat
Step-1 run failed the same way and is likewise excluded.

Finally, a note on the report's status. Two defects in the measurement apparatus were
identified during this analysis and have since been corrected in the pipeline: the
prompt field order was uncontrolled (Section~\ref{sec:ordering}), and the LLM client
declared no token cap (Section~\ref{sec:maxtokens}). Every result reported here
predates both corrections. The report is therefore a complete account of what the
current evidence supports, together with an explicit statement of which conclusions
those corrections put back in question; Section~\ref{sec:status} sets out what
remains to be run.

\section{Pipeline Overview}

The pipeline runs in five stages, each timed and recorded in a per-run JSON summary
under \texttt{logs/<dataset>/runs/}.

\textbf{Load and normalize.} Both tables and the ground truth are read, ID columns
are dropped, and every remaining field is passed through a normalization step that
expands or standardizes address and directional abbreviations, and cleans up
Unicode/HTML artifacts, so that superficial formatting differences do not by
themselves create false mismatches later in the pipeline.

\textbf{Embed and block.} Each attribute column is encoded independently with the
SentenceBERT model, the column vectors are concatenated and L2-normalized into one
embedding per record, and multi-probe locality-sensitive hashing (random
hyperplanes, with single-bit flips at query time) retrieves the top-$k$ most similar
Table-B candidates for every Table-A record. This reduces the pairwise comparison
space from the full cross product down to a much smaller candidate-pair set before
any LLM is involved.

\textbf{Step 1 -- global tuple condensation (optional).} One attribute subset is
chosen once for the whole dataset and applied to every candidate pair before
re-blocking on the condensed attributes. This step is described in full in
Section~\ref{sec:step1}.

\textbf{Step 1.5 -- adaptive per-pair attribute selection (optional).} Instead of one
global subset, a small sample of candidate pairs is profiled by the LLM for
per-attribute importance, and a hybrid predictor/retriever model transfers a
\emph{different} attribute subset to every candidate pair. This step is described in
full in Section~\ref{sec:step15}.

\textbf{Match.} One LLM call decides Yes/No per candidate pair, using only the
attributes selected for that pair (all attributes if neither Step 1 nor Step 1.5
ran). Every response is cached on disk keyed by a hash of the model, prompt
template, and record content, so repeated experiments over the same configuration
do not re-spend tokens.

\textbf{Metrics and logging.} Precision, recall, and F1 are computed both at the
candidate-pair level (only pairs that survived blocking) and end-to-end (relative to
all ground-truth matches, so blocking misses count against recall), alongside
token-usage accounting split by stage.

\section{Methodology}
\label{sec:methodology}

\subsection{Step 1: Global Tuple Condensation}
\label{sec:step1}

The motivation for Step 1 is that showing an LLM every attribute for every pair
wastes tokens on fields that are frequently empty, redundant, or simply irrelevant
to the match decision, and can make the model's decision worse by diluting the
signal from the few attributes that actually distinguish entities. Step 1 picks
\emph{one} attribute subset for the whole dataset, used for every pair, via one of
four strategies.

\paragraph{Manual selection.} A fixed, user-supplied list of attribute names is kept
and everything else is dropped, with no data-driven reasoning involved. It is the
cheapest strategy to run (no LLM calls, no labeled data required) and is appropriate
when a person can already tell from the schema which fields carry identity
information, but it does not adapt if that judgment is wrong for some records.

\paragraph{Heuristic selection.} Attributes are pruned using two label-free,
corpus-level statistics: the fraction of missing values (an attribute that is empty
in most records cannot be evidence of anything) and the fraction of unique values
(an attribute that takes almost the same value everywhere carries little
discriminative signal). Attributes that survive pruning are also \emph{compressed}:
long text fields are stripped of boilerplate words (via a stopword/marketing-language
regex list) and only high-signal tokens are kept -- numbers, alphanumeric
codes/identifiers, capitalized proper nouns, and unit-bearing measurements (e.g.\
\texttt{"12 kg"}) -- so that a lengthy description field still contributes its
distinctive tokens without its full token count.

\paragraph{LLM-guided selection.} A small labeled sample of matching and
non-matching pairs is shown to the LLM one pair at a time, together with the
question ``which fields were most influential in deciding whether these two records
match?'', and the model returns the smallest sufficient subset of attribute names for
that pair. Responses are aggregated across the sample by \emph{selection frequency}
per attribute (what fraction of responses named that attribute), attributes are
ranked by that frequency, and the final subset keeps every attribute above a
frequency threshold (optionally capped to a top-$k$).

\paragraph{Supervised selection.} For the same labeled sample, three field-level
similarity features are computed per attribute for every pair -- token-set Jaccard
overlap, normalized edit-distance similarity, and an exact-match indicator -- and a
logistic regression classifier (features standardized, class-balanced) is trained to
predict match/non-match from all attributes' features at once. Each attribute's
importance is the maximum absolute regression weight across its three features,
normalized so all attributes' importance scores sum to one; the top-ranked
attributes above a threshold (optionally capped to top-$k$) are kept. Because it
learns from labeled evidence directly, this strategy tends to identify a very small,
highly discriminative attribute subset (see Section~\ref{sec:results}).

Whichever strategy is used, the condensed table is then re-embedded and re-blocked
(so blocking itself also benefits from the smaller, more focused attribute set)
before matching.

\subsection{Step 1.5: Pair-Adaptive Attribute Selection}
\label{sec:step15}

Step 1.5 addresses a limitation of Step 1: even a well-chosen global attribute
subset is not equally informative for every pair. For one pair an exact model
number may already be decisive; for another, the deciding signal might be an author
mismatch or a subtle title variation, while the model number field is missing or
uninformative. Step 1.5 replaces one global subset with a \emph{per-pair} subset,
produced in two stages.

\textbf{Profiling stage.} A small sample of post-blocking candidate pairs (the
``profile pairs'') is selected \emph{without using ground-truth labels} by one of
five samplers (described below). For each profile pair, the LLM is asked to score
every attribute's usefulness for the match/non-match decision on an integer 0--3
scale (0 = not useful/missing/generic, 3 = highly important or decisive), and the
raw scores are normalized per pair into an importance vector that sums to one.

\textbf{Transfer stage.} A \texttt{HybridAttributeSelector} is fit on the profiled
pairs' label-free similarity features (token overlap, edit similarity, exact match,
semantic similarity, numeric similarity, and an is-numeric indicator, computed per
attribute) against their LLM-scored importance vectors. It combines two
sub-models: a Ridge-regression \emph{predictor} that maps a pair's features directly
to a predicted importance vector, and a cosine-similarity \emph{retriever} that finds
the $k$ most similar profiled pairs and averages their importance vectors. The two
signals are combined by a weighted average (0.5/0.5 in the current configuration).
For every candidate pair (profiled or not), attributes are then selected in
descending combined-importance order until either a fixed top-$k$ is reached or
cumulative importance crosses a threshold (0.8 in the current configuration,
bounded between a minimum and maximum attribute count), with a fallback rule that
force-includes at least one attribute tagged with a required ``identity'' role
(e.g.\ title, name, model number) even if the learned ranking would otherwise drop
it.

The five samplers below differ only in \emph{which} candidate pairs get profiled by
the LLM; the transfer stage is identical for all of them.

\paragraph{Random sampling.} Profile pairs are drawn uniformly at random from all
candidate pairs. This is the label-free baseline against which the other samplers
are compared.

\paragraph{Similarity-stratified sampling.} Candidate pairs are bucketed into
low/medium/high terciles by their blocking-stage cosine similarity, and profile
pairs are drawn from each bucket in fixed proportions (60\% medium, 20\% low, 20\%
high in the current configuration). This intentionally profiles some clearly
similar pairs, some clearly dissimilar pairs, and a majority of ambiguous
near-threshold pairs, on the reasoning that ambiguous pairs are where attribute
importance is hardest to guess and most useful to learn from.

\paragraph{Diversity sampling.} Profile pairs are chosen by farthest-first traversal
in the label-free pair-feature space: starting from the pair farthest from the
sample centroid, each subsequent pair is the one that maximizes its minimum
distance to all previously chosen pairs. This favors broad coverage of distinct
similarity ``shapes'' (e.g.\ pairs that agree on names but disagree on numbers, versus
pairs that disagree on everything) over any particular similarity level.

\paragraph{Stratified diversity sampling.} The same low/medium/high similarity
strata as similarity-stratified sampling are used, but within each stratum, pairs
are chosen by the same farthest-first diversity rule rather than uniformly at
random. This is meant to combine both benefits: similarity-level balance and
within-level feature diversity.

\paragraph{Importance-based sampling.} Pairs are ranked by a label-free
\emph{importance proxy score} that estimates, without ground truth, how
informative a pair is likely to be for learning attribute importance:
$$\text{proxy} = 0.50 \cdot \text{evidence\_contrast} + 0.30 \cdot \text{similarity\_uncertainty} + 0.20 \cdot \text{evidence\_mixedness}$$
%
where evidence\_contrast is the standard deviation of per-attribute agreement
signals (favoring pairs where attributes disagree with each other, rather than all
pointing the same way), similarity\_uncertainty favors pairs whose blocking
similarity is far from the sample's typical value, and evidence\_mixedness favors
pairs whose average per-attribute agreement sits near 0.5 (neither clearly matching
nor clearly non-matching). The top-scoring pairs (with random tie-breaking) become
the profile sample.

\section{Experimental Setup}
\label{sec:setup}

Table~\ref{tab:config} lists the current per-dataset configuration used for every
number in this report. Table~\ref{tab:provenance} lists the exact run directory each
result was read from, for traceability.

\begin{table}[htbp]
\centering
\small
\begin{tabular}{lrrrrr}
\toprule
Dataset & Blocking (tables/planes/flips/top-$k$) & Profile size & Attrs. available \\
\midrule
Fodors--Zagat   & 5 / 8 / 1 / 5  & 20 & 5 (6 less \texttt{class}) \\
DBLP--ACM       & 15 / 8 / 1 / 5 & 50 & 4 \\
Amazon--Walmart & 10 / 4 / 1 / 4 & 50 & 13 \\
\bottomrule
\end{tabular}
\caption{Current per-dataset blocking and Step-1.5 profiling configuration.
Adaptive selection policy is shared across datasets: top-$k$ retrieval~5, min/max
attributes 2/5, cumulative-importance threshold~0.8, required role~identity.}
\label{tab:config}
\end{table}

\begin{table}[htbp]
\centering
\small
\resizebox{\textwidth}{!}{%
\begin{tabular}{lll}
\toprule
Dataset & Method & Run directory \\
\midrule
DBLP--ACM & step15 / random & \texttt{20260615\_120346\_step15\_random} \\
DBLP--ACM & step15 / similarity\_stratified & \texttt{20260615\_121510\_step15\_similarity\_stratified} \\
DBLP--ACM & step15 / diversity & \texttt{20260615\_123251\_step15\_diversity} \\
DBLP--ACM & step15 / stratified\_diversity & \texttt{20260615\_133158\_step15\_stratified\_diversity} \\
DBLP--ACM & step15 / importance\_based & \texttt{20260615\_221942\_step15\_importance\_based} \\
DBLP--ACM & step1 / supervised & \texttt{20260615\_225756\_step1} \\
DBLP--ACM & step15 / random, 5 seeds & \texttt{20260802\_035813}--\texttt{20260802\_044306} (5 dirs) \\
DBLP--ACM & \quad \emph{aborted, API exhaustion} & \texttt{20260802\_045249}--\texttt{20260802\_051239} (8 dirs) \\
Fodors--Zagat & step15 / all 5 samplers, 6 runs each & \texttt{20260801\_2246*}--\texttt{20260802\_0032*} (30 dirs) \\
Fodors--Zagat & \quad aggregated in & \texttt{logs/Fodors-Zagat/sampler\_seed\_sweep.json} \\
Fodors--Zagat & step1 / supervised & \texttt{20260801\_230544\_step1} \\
Fodors--Zagat & full / no selection & \texttt{20260801\_230923\_full} \\
Fodors--Zagat & full / no selection, replicate & \texttt{20260802\_173229\_full} \\
Amazon--Walmart & step15 / similarity\_stratified & \texttt{20260615\_163220\_step15\_similarity\_stratified} \\
\bottomrule
\end{tabular}%
}
\caption{Provenance: every metric in this report is read from the \texttt{run\_summary.json} in the listed \texttt{logs/<dataset>/runs/} directory. All Fodors--Zagat runs listed were executed with \texttt{--exclude-attributes class}; all runs listed were executed with \texttt{--disable-matcher-cache}. The eight aborted DBLP--ACM directories are recorded for completeness but contribute no metrics; see Section~\ref{sec:dblp-seeds}.}
\label{tab:provenance}
\end{table}

\section{Results}
\label{sec:results}

Table~\ref{tab:step15-results} compares the five Step-1.5 profile samplers on
DBLP--ACM. Blocking pair completeness is 1.0000 for every row shown (every
ground-truth match survived blocking), so candidate-pair metrics and end-to-end
metrics are identical here. Bold marks the best value per column.

\begin{table}[htbp]
\centering
\small
\resizebox{\textwidth}{!}{%
\begin{tabular}{llrrrrrrr}
\toprule
Dataset & Sampler & Attrs & F1 & Precision & Recall & TP & FP & FN \\
\midrule
DBLP--ACM & random & 4 & 0.9235 & 0.9722 & 0.8795 & 1{,}956 & 56 & 268 \\
DBLP--ACM & similarity\_stratified & 4 & 0.9200 & 0.9715 & 0.8737 & 1{,}943 & 57 & 281 \\
DBLP--ACM & diversity & 4 & 0.9162 & 0.9761 & 0.8633 & 1{,}920 & \textbf{47} & 304 \\
DBLP--ACM & stratified\_diversity & 4 & 0.9066 & 0.9669 & 0.8534 & 1{,}898 & 65 & 326 \\
DBLP--ACM & importance\_based & 4 & \textbf{0.9605} & \textbf{0.9768} & \textbf{0.9447} & \textbf{2{,}101} & 50 & \textbf{123} \\
\bottomrule
\end{tabular}%
}
\caption{Step-1.5 profile-sampler comparison on DBLP--ACM, current pipeline
(\texttt{BAAI/bge-large-en-v1.5} + \texttt{google/gemini-2.5-flash}). Single run per
sampler at seed 42. Section~\ref{sec:dblp-seeds} shows that this single-seed ranking
is not reliable.}
\label{tab:step15-results}
\end{table}

\subsection{Reseeding DBLP--ACM: A Partial Sweep}
\label{sec:dblp-seeds}

A repeated-seed sweep was attempted on DBLP--ACM to test whether the ranking in
Table~\ref{tab:step15-results} survives reseeding. It completed only in part. The
five \texttt{random} runs (seeds 42--46) finished cleanly with zero API errors. The
sixth run, \texttt{similarity\_stratified} at seed 42, failed on 3{,}785 of its
13{,}080 pairs (28.9\%), and the remaining eight runs failed on all 13{,}080 pairs,
consistent with the OpenRouter budget being exhausted mid-sweep. The pipeline scores
a failed LLM call as a non-match and still writes a complete
\texttt{run\_summary.json}, so those nine runs carry plausible-looking but
meaningless metrics (the fully-failed runs report F1 $=0.0000$; the partial run
reports $0.7672$). They are excluded throughout. Only the \texttt{random} rows in
Table~\ref{tab:dblp-seeds} are evidence.

\begin{table}[htbp]
\centering
\small
\resizebox{\textwidth}{!}{%
\begin{tabular}{lrrrrr}
\toprule
Seed & F1 & Precision & Recall & TP & FN \\
\midrule
42 & 0.9258 & 0.9718 & 0.8840 & 1{,}966 & 258 \\
43 & 0.9170 & 0.9766 & 0.8642 & 1{,}922 & 302 \\
44 & 0.8847 & 0.9710 & 0.8125 & 1{,}807 & 417 \\
45 & 0.9150 & 0.9760 & 0.8611 & 1{,}915 & 309 \\
46 & 0.9119 & 0.9749 & 0.8566 & 1{,}905 & 319 \\
\midrule
mean & 0.9109 & 0.9741 & 0.8557 & 1{,}903 & 321 \\
SD & 0.0155 & 0.0025 & 0.0264 & 59 & 59 \\
range & 0.0411 & 0.0056 & 0.0715 & 159 & 159 \\
\bottomrule
\end{tabular}%
}
\caption{DBLP--ACM, \texttt{random} profile sampler over five seeds. Blocking is
held fixed across all five runs, so the candidate-pair set is identical and the only
varying input is which 50 pairs are profiled.}
\label{tab:dblp-seeds}
\end{table}

The spread is large. Reseeding the profile sample alone moves end-to-end F1 by
0.0411 and true positives by 159 --- comfortably wider than the 0.0370--0.0539 gap
that separates \texttt{importance\_based} from the other four samplers in
Table~\ref{tab:step15-results}. A single-seed sampler ranking on this dataset is
therefore not interpretable.

This is not LLM non-determinism. The June and August \texttt{random} runs at seed 42
share identical configuration and produce identical attribute selections on 99.9\% of
pairs; their answers differ on 0.22\% of pairs, worth 10 true positives and 0.0023
F1. Replicate noise is roughly one-eighteenth of the seed-to-seed spread.

\subsection{Prompt Attribute Order Is an Uncontrolled Variable}
\label{sec:ordering}

Decomposing the seed-to-seed disagreements locates the cause, and it is not the one
the experiment was designed to measure. \texttt{HybridAttributeSelector.select\_attributes}
ranks attributes by predicted importance (\texttt{np.argsort(-combined)}) and emits
the selected subset \emph{in that order}. \texttt{code/matcher.py:116} rebuilds each
record as a dict comprehension over that list and serializes it with
\texttt{json.dumps}, so the ranking order is reproduced verbatim as the field order
in the prompt. Because predicted importance shifts continuously with the profile
sample, reseeding permutes the prompt even when it selects exactly the same
attributes.

Table~\ref{tab:order-effect} conditions every cross-seed pair comparison on whether
the two runs selected the same attribute \emph{set} and the same \emph{order}.

\begin{table}[htbp]
\centering
\small
\resizebox{\textwidth}{!}{%
\begin{tabular}{lrr}
\toprule
Condition & Mean pairs per seed comparison & Answers disagree \\
\midrule
Same attribute set, same order & 1{,}871 & \textbf{0.10\%} \\
Same attribute set, different order & 4{,}772 & \textbf{5.44\%} \\
Different attribute set & 6{,}437 & 1.23\% \\
\bottomrule
\end{tabular}%
}
\caption{DBLP--ACM, \texttt{random} sampler, all ten pairwise comparisons among the
five seeds (13{,}080 candidate pairs each). Holding the selected attribute set fixed
and permuting only the order in which those attributes appear in the prompt changes
the matcher's answer 54 times more often than presenting the identical prompt twice.}
\label{tab:order-effect}
\end{table}

Presenting the identical prompt twice flips 0.10\% of answers; presenting the same
attributes in a different order flips 5.44\%. Changing \emph{which} attributes are
selected flips 1.23\% --- less than reordering does, though that group is not a
controlled comparison, since the pairs on which two selectors disagree are also the
pairs where they were least confident. The controlled contrast is the first two rows,
and it is unambiguous: on this schema the matcher is more sensitive to field order
than to field choice.

Order also varies far more than set does. Across seed comparisons the two runs agree
on the attribute set for 42--67\% of pairs, but among those agreements the order
still differs on 51--96\%. DBLP--ACM has only four attributes, and the selector is in
practice making a single binary decision per pair --- keep \texttt{year} or not
(mean selection size 3.51--3.79, never below 3 or above 4). The remaining variation
between runs is permutation.

The consequence is that Step-1.5's measured effect on DBLP--ACM is confounded. Each
run varies two things at once: the attribute subset, which is the intended
independent variable, and the prompt field order, which is an unintended by-product
of ranking and which by this measurement is the stronger of the two. The same
confound is present in every Step-1.5 number in this report, Fodors--Zagat included,
and has not been quantified there.

This has since been corrected in the pipeline. Both selection stages now sort their
chosen subset into table-column order before condensation, so field order is fixed
across runs and attribute choice is the only thing a reseed varies; the
importance ranking is still recorded, but no longer reaches the prompt. Runs produced
after the fix carry \texttt{canonical\_attribute\_order} in their run summary, which
distinguishes the two generations. Every result in this report predates the fix, so
the confound stands as described above; re-running the comparisons is the first item
of remaining work (Section~\ref{sec:status}).

\subsection{An Unset Token Cap Made Runs Fail on Budget}
\label{sec:maxtokens}

A second infrastructure defect surfaced while attempting the re-runs and is recorded
here because it explains the aborted sweeps of Section~\ref{sec:dblp-seeds}. The
shared OpenRouter client never set \texttt{max\_tokens}, so each request defaulted to
the model's full output ceiling of 65{,}535 tokens. OpenRouter reserves credit against
that declared ceiling rather than against actual usage, and rejects a request with
HTTP 402 whenever the remaining balance cannot cover it.

The reservation was pathologically conservative. Measured over clean runs, the matcher
emits \emph{exactly one} completion token on every call --- 13{,}080 of 13{,}080 on
DBLP--ACM, 2{,}665 of 2{,}665 on Fodors--Zagat --- because it answers with a single
label. Attribute-importance profiling emits 50 tokens on DBLP--ACM's four-attribute
schema and 133 on Amazon--Walmart's thirteen. Every call therefore reserved roughly
four orders of magnitude more budget than it consumed, and runs began failing while
the account still held usable credit.

The client now sends an explicit cap: 8 tokens for matching, 512 for profiling. A
post-fix \texttt{--mode full} run on Fodors--Zagat completed with zero API errors and
2{,}665 completion tokens over 2{,}665 calls, confirming the cap is not truncating
responses. This changes no experimental result --- the caps sit far above observed
usage --- but it removes a failure mode that silently corrupted nine runs.

Fodors--Zagat was run six times per sampler (seeds 42--46, plus one repeat of seed
42), so it is reported as a distribution rather than a point estimate in
Table~\ref{tab:fodors-seeds}. This matters: on a 110-pair ground truth a single true
positive is worth 0.9 percentage points of recall, and the single-seed figures that
this report previously carried turned out not to be representative for two of the
five samplers.

\begin{table}[htbp]
\centering
\small
\resizebox{\textwidth}{!}{%
\begin{tabular}{lrrrrrrr}
\toprule
Sampler & Mean F1 & SD & Min & Max & TP range & Mean prec. & Mean rec. \\
\midrule
diversity              & \textbf{0.9222} & \textbf{0.0021} & \textbf{0.9179} & 0.9231 & 95--96 & 0.9796 & \textbf{0.8712} \\
stratified\_diversity  & 0.9152 & 0.0042 & 0.9135 & 0.9238 & 95--97 & 0.9695 & 0.8667 \\
similarity\_stratified & 0.9110 & 0.0211 & 0.8788 & \textbf{0.9384} & 87--99 & 0.9827 & 0.8500 \\
random                 & 0.8843 & 0.0429 & 0.8316 & 0.9275 & 79--96 & 0.9854 & 0.8045 \\
importance\_based      & 0.8356 & 0.0063 & 0.8254 & 0.8437 & 78--81 & \textbf{0.9876} & 0.7242 \\
\bottomrule
\end{tabular}%
}
\caption{Fodors--Zagat (\texttt{class} excluded), Step-1.5 profile samplers over six
runs each. End-to-end metrics; blocking pair completeness is 1.0000 throughout, and
blocking is held fixed across seeds so that only the profiling sample varies.}
\label{tab:fodors-seeds}
\end{table}

Amazon--Walmart only has one Step-1.5 run on the current pipeline
(similarity\_stratified), so it is reported separately rather than alongside the
five-way comparison above. Pair completeness is 0.9099, so 103 ground-truth matches
were already lost during blocking, before the LLM ever saw them; candidate-level and
end-to-end metrics therefore differ.

\begin{table}[htbp]
\centering
\small
\begin{tabular}{lrrrrrr}
\toprule
Metric level & Attrs & Precision & Recall & F1 & TP & FN \\
\midrule
Candidate-pair & 13 & 0.8588 & 0.8596 & 0.8592 & 894 & 146 \\
End-to-end     & 13 & 0.8588 & 0.7822 & 0.8187 & 894 & 249 \\
\bottomrule
\end{tabular}
\caption{Amazon--Walmart, Step-1.5 similarity\_stratified sampler (only current-pipeline run available). End-to-end FN includes the 103 ground-truth matches lost at blocking.}
\label{tab:amazon-results}
\end{table}

Table~\ref{tab:step1-vs-step15} compares Step 1's supervised global condensation
against the best Step-1.5 sampler on each of the two datasets where both a current
Step 1 and a current Step 1.5 sweep exist.

\begin{table}[htbp]
\centering
\small
\resizebox{\textwidth}{!}{%
\begin{tabular}{llrrrrr}
\toprule
Dataset & Method & Attrs used & F1 & Precision & Recall & Tokens \\
\midrule
DBLP--ACM & Step 1, supervised (authors, title) & 2 & 0.9515 & 0.9289 & \textbf{0.9753} & \textbf{3{,}502{,}759} \\
DBLP--ACM & Step 1.5, importance\_based (best) & 4 & \textbf{0.9605} & \textbf{0.9768} & 0.9447 & 3{,}910{,}646 \\
\midrule
Fodors--Zagat & Step 1, supervised (phone, name) & 2 & \textbf{0.9537} & 0.9717 & \textbf{0.9364} & \textbf{658{,}185} \\
Fodors--Zagat & Step 1.5, diversity (best mean of 6 runs) & 5 & 0.9222 & \textbf{0.9796} & 0.8712 & 775{,}828 \\
\bottomrule
\end{tabular}%
}
\caption{Global (Step 1) versus adaptive per-pair (Step 1.5, best sampler) attribute selection. F1/precision/recall are end-to-end.}
\label{tab:step1-vs-step15}
\end{table}

Fodors--Zagat is the only dataset with a no-selection baseline
(\texttt{--mode full}), which makes it the only place where the absolute value of
doing attribute selection at all can be measured. Table~\ref{tab:fodors-baseline}
places that baseline alongside the two selection mechanisms.

\begin{table}[htbp]
\centering
\small
\resizebox{\textwidth}{!}{%
\begin{tabular}{lrrrrrr}
\toprule
Method & Attrs/pair & F1 & Precision & Recall & Tokens & vs.\ no selection \\
\midrule
No selection (\texttt{full}) & 5.00 & 0.7849 & 0.9605 & 0.6636 & 802{,}477 & --- \\
Step 1.5, diversity (mean of 6) & 4.00 & 0.9222 & 0.9796 & 0.8712 & 775{,}828 & $+0.1373$ F1, $-3.3\%$ tokens \\
Step 1, supervised & 2.00 & \textbf{0.9537} & 0.9717 & \textbf{0.9364} & \textbf{658{,}185} & $+0.1688$ F1, $-18.0\%$ tokens \\
\bottomrule
\end{tabular}%
}
\caption{Fodors--Zagat (\texttt{class} excluded): attribute selection versus matching on all attributes. F1/precision/recall are end-to-end. ``Attrs/pair'' is the mean number of attributes shown to the matcher LLM per candidate pair.}
\label{tab:fodors-baseline}
\end{table}

The no-selection baseline has since been run a second time under an identical
configuration, giving an independent estimate of the matcher's replicate noise on this
dataset. The two runs return F1 of 0.7849 and 0.7826 (73 and 72 true positives, 3 and
2 false positives) on identical token spend. Because \texttt{--mode full} applies no
selection, the canonical-order fix of Section~\ref{sec:ordering} cannot affect it, and
the 0.0023 gap is attributable entirely to LLM non-determinism. One true positive out
of 110 is consistent with the 0.10--0.22\% per-pair disagreement measured on
DBLP--ACM, and it sets the scale against which the Fodors--Zagat sampler differences
in Table~\ref{tab:fodors-seeds} should be read.

\subsection{Step 1 and Step 1.5 Do Not Operate Under the Same Supervision}
\label{sec:supervision}

Every head-to-head between the two mechanisms in this report compares a method that
consumes ground-truth labels against one that does not, and the gap in supervision is
wide enough to change what the comparison establishes.

Step 1's supervised strategy calls \texttt{build\_step1\_labeled\_pairs}
(\texttt{code/main.py:144}), which draws its positive examples directly from the gold
set and its negatives from blocked non-matches, then trains a logistic regression on
the result. The LLM-guided strategy consumes the same labeled sample. Step 1.5 uses no
labels anywhere: profile pairs are chosen by label-free samplers over label-free pair
features, and are scored by the LLM rather than against ground truth.

\begin{table}[htbp]
\centering
\small
\begin{tabular}{lrrr}
\toprule
Mechanism & Gold labels consumed & Fodors--Zagat & DBLP--ACM \\
\midrule
Step 1, supervised / LLM-guided & 50 pos.\ $+$ 50 neg. & 45.5\% of gold & 2.2\% of gold \\
Step 1, manual / heuristic & none & --- & --- \\
Step 1.5, all five samplers & none & --- & --- \\
\bottomrule
\end{tabular}
\caption{Ground-truth consumption by attribute-selection mechanism. Percentages are
the 50 sampled positives as a share of each dataset's gold matches (110 and 2{,}224
respectively).}
\label{tab:supervision}
\end{table}

On Fodors--Zagat the asymmetry is severe. The ground truth contains 110 matches and
the supervised strategy is fitted on 50 of them --- it sees 45.5\% of the answer key
before selecting an attribute. On DBLP--ACM the same 50 positives are only 2.2\% of a
much larger gold set, and the concern is correspondingly milder.

There is a second issue beyond the asymmetry. The labeled sample is not held out:
\texttt{build\_step1\_labeled\_pairs} draws from \texttt{gt\_set}, and the evaluation
at \texttt{code/main.py:1032} scores against that same \texttt{gt\_set} without
excluding the sampled pairs. Part of the Step-1 score is therefore measured on pairs
whose attribute selection was fitted with knowledge of their labels. This is not
catastrophic --- what is fitted is a global attribute ranking, not a per-pair decision,
so the leak is indirect --- but on Fodors--Zagat it touches 45.5\% of the evaluation
set and should be stated rather than left implicit.

Neither point makes the Step-1 numbers wrong, and neither rescues Step 1.5, which must
still demonstrate a win on its own terms. What they change is the question. Step 1.5
is not competing with a peer; it is competing with a method handed a large share of
the ground truth, which makes Step 1 supervised closer to an approximate upper bound
than to a deployable baseline --- in any real deployment the labels it consumes are
precisely what the pipeline exists to produce. Read that way, 0.9222 against 0.9537 on
Fodors--Zagat under zero supervision is a materially different result from the same
gap between two comparable methods.

The genuine like-for-like baselines are Step 1's \emph{manual} and \emph{heuristic}
strategies, both label-free. Neither has a current-pipeline run on any dataset
(Section~\ref{sec:limitations}), so the evaluation as it stands cannot answer the
question the thesis most needs answered.

\subsection{Label Leakage in Fodors--Zagat}
\label{sec:leakage}

The Fodors--Zagat schema contains an attribute named \texttt{class} which, on
inspection, is the ground-truth cluster identifier rather than a data field:
\texttt{class} takes the same value on both sides for 110 of 110 ground-truth pairs
and is unique within each table. Matching on \texttt{class} alone therefore solves
the dataset exactly, and any method that selects it is being evaluated on a label
rather than on evidence.

The Step-1 supervised strategy selected it. Its ranking assigned \texttt{phone} and
\texttt{class} the identical importance score 0.2825, the expected signature of two
perfectly predictive features, and the resulting run reported 0.9725 F1 -- the
highest figure of any method on any dataset in the first generation of runs. That
number is not a valid measurement of the method.

The Step-1.5 path was affected differently. The profiling LLM consistently rated
\texttt{class} the \emph{least} useful attribute of the six (mean importance 1.35
and 1.50 out of 3 across the two runs, against 3.00 for \texttt{phone}), so it
declined to exploit the leak. The leak nonetheless degraded Step 1.5 by occupying a
selection slot with a field the matcher could not interpret: after removing
\texttt{class}, four of five samplers improved substantially, as shown in
Table~\ref{tab:leak-effect}.

\begin{table}[htbp]
\centering
\small
\begin{tabular}{lrrr}
\toprule
Method & F1 with \texttt{class} & F1 without \texttt{class} & $\Delta$ \\
\midrule
Step 1.5, random & 0.7802 & 0.9179 & $+0.1377$ \\
Step 1.5, similarity\_stratified & 0.9171 & 0.9231 & $+0.0060$ \\
Step 1.5, diversity & 0.8254 & 0.9231 & $+0.0977$ \\
Step 1.5, stratified\_diversity & 0.8108 & 0.9135 & $+0.1027$ \\
Step 1.5, importance\_based & 0.8601 & 0.8377 & $-0.0224$ \\
Step 1, supervised & 0.9725$^\dagger$ & 0.9581 & $-0.0144$ \\
\bottomrule
\end{tabular}
\caption{Effect of removing the leaking \texttt{class} attribute on Fodors--Zagat.
F1 is candidate-level for comparability. $^\dagger$This figure is contaminated:
the supervised strategy selected \texttt{class} itself. With \texttt{class} removed
it selects (\texttt{phone}, \texttt{name}) instead.}
\label{tab:leak-effect}
\end{table}

Two consequences follow. First, the first-generation Fodors--Zagat sampler ranking
was largely an artifact: removing the leak changes which sampler wins and compresses
the spread between them from 0.137 F1 to 0.085. Second, the qualitative conclusion
that global condensation outperforms adaptive selection on this dataset survives the
correction -- Step 1 still leads, but by 0.031 F1 rather than 0.055, and now on a
defensible attribute pair.

\section{Interpretation}

\subsection{Fodors--Zagat}

The six-run sweep separates two effects that a single run confounds: how accurate a
sampler is on average, and how much its accuracy depends on which pairs it happens
to draw.

On the first, only one difference is unambiguous. Importance-based sampling is
significantly worse than all four other samplers (Welch $t$ from $-2.75$ against
random to $-25.63$ against stratified diversity), and it is worse by a wide margin:
its best run of six (0.8437) falls below the worst run of every sampler except
random. The remaining four are statistically indistinguishable from each other in
mean, with one marginal exception (diversity over stratified diversity, $t = 3.64$,
but a difference of only 0.0070 F1). The apparent ordering among those four in any
single run is not a finding.

On the second, the differences are large and systematic, and they track exactly
which samplers are stochastic. Repeating a run at a fixed seed and configuration
changes the result by 0--3 true positives, which bounds the matcher LLM's own
non-determinism. Against that floor, the deterministic samplers barely move across
seeds (diversity SD 0.0021, stratified diversity 0.0042, importance-based 0.0063),
whereas the two that genuinely resample vary by an order of magnitude more
(similarity-stratified SD 0.0211, random SD 0.0429). A single run of random
sampling can land anywhere between 79 and 96 true positives -- a span that covers
almost the entire distance between the best and worst samplers. Random and
similarity-stratified sampling are therefore not merely less accurate on average;
they are unreliable when only one run is affordable, which is the operating regime
Step 1.5 is designed for.

This reframes what the sampler comparison contributes. The useful property of
diversity and stratified-diversity sampling is not a higher mean but a much lower
variance: by choosing profile pairs deterministically from the feature geometry
rather than by drawing them, they remove the dominant source of run-to-run
instability. That the single best individual run in the sweep belongs to
similarity-stratified sampling (0.9384) does not contradict this; it is what a
high-variance sampler looks like when a draw goes well.

The mean differences that do exist reduce to a single binary decision. Every sampler
selects \texttt{name}, \texttt{addr}, and \texttt{phone} for essentially every
pair; the only variation is whether the fourth slot goes to \texttt{city} or
\texttt{type}. On the 110 true matches, random, diversity, and stratified-diversity
choose \texttt{city} for 96--100\% of pairs, similarity-stratified for 86\%, and
importance-based for only 22\%. That substitution is what separates the samplers:
across the five seed-42 runs, true matches shown (\texttt{addr}, \texttt{name},
\texttt{phone}, \texttt{city}) are missed 9.5--12.7\% of the time, whereas true
matches shown (\texttt{addr}, \texttt{name}, \texttt{phone}, \texttt{type}) are
missed 29--67\% of the time. \texttt{type} is a low-cardinality cuisine category
whose taxonomy differs between the two sources, so it disagrees on genuine matches
and pushes the matcher toward ``No''.

Importance-based sampling is the clear loser here (mean 0.8356 F1 over six runs, 29
to 32 false negatives against 13--15 for the others) purely because it makes that
substitution most often, and it does so consistently: because its profile selection
is deterministic, the poor choice is reproduced on every seed rather than being a
bad draw. Notably, this is not a transfer-stage failure in the mechanical sense:
leave-one-out diagnostics on its own profile set give it the \emph{highest} internal
fidelity of any Fodors--Zagat run (rank correlation 0.82 and selected-set Jaccard
0.96 between predicted and LLM-assigned importance). The transfer stage reproduces
its profile sample faithfully; the profile sample is simply unrepresentative,
because the proxy score deliberately selects high-contrast, contested pairs, and on
those pairs \texttt{type} does carry apparent signal.

Against the no-selection baseline the case for doing selection at all is strong.
Matching on all five attributes yields only 0.7849 F1 (recall 0.6636); every
selection method beats it by 0.13--0.17 F1, and even the single worst run in the
entire six-seed sweep (0.8316) beats it. The token picture is less flattering to
Step 1.5: dropping one of five attributes per pair saves only 3.3\% of tokens,
because the attribute it drops (\texttt{city} or \texttt{type}) is among the
shortest fields. Step 1's two-attribute global subset saves 18.0\% by comparison.
Attribute count and token count are not proportional, and Step 1.5's current
selection policy optimizes the former.

\subsection{DBLP--ACM}

At single seed, DBLP--ACM appeared to be the dataset where Step 1.5 worked best:
importance-based sampling reached 0.9605 F1 against 0.90--0.92 for the other four
samplers, driven almost entirely by recall (0.9447 versus 0.85--0.88), and edged out
Step 1's supervised condensation (0.9515 F1 on authors + title alone) at the cost of
roughly 12\% more matching tokens. That reading no longer holds up.

The reseeding sweep (Section~\ref{sec:dblp-seeds}) shows that changing the profile
sample alone moves \texttt{random}'s F1 across 0.8847--0.9258, a 0.0411 range. The
gap the single-seed table attributes to sampler choice is 0.0370--0.0539. The
measurement is not sharp enough to resolve the effect it was reporting, so the
headline that importance-based sampling wins on DBLP--ACM must be withdrawn pending
repeated seeds --- not contradicted, but unsupported. It may well survive: it is a
deterministic sampler, and the one comparison available at fixed configuration
(June versus August \texttt{random} at seed 42) puts replicate noise at only 10 true
positives. But that has not been shown.

Two structural facts limit what Step 1.5 can express here in any case. First, the
schema has four attributes and the policy floor is \texttt{min\_k}${}=2$ with
\texttt{title} and \texttt{authors} selected essentially always, so the selector's
entire per-pair decision reduces to whether to include \texttt{year}: observed
selection sizes are 3 or 4 and nothing else. Second, \texttt{year} agrees on 100\% of
the 2{,}224 ground-truth matches, making it a low-information field --- including or
excluding it should not, and largely does not, change the answer. The apparent
recall differences between the $k=3$ and $k=4$ groups within a run (0.84--0.98 versus
0.81--0.90) reverse direction between seeds and reflect which pairs the selector
routed where, not a causal effect of dropping \texttt{year}.

What does move the answers is prompt field order (Section~\ref{sec:ordering}). On a
four-attribute schema where the attribute set is nearly fixed, permutation is most of
what reseeding actually varies, and it flips 5.44\% of answers. The DBLP--ACM
Step-1.5 results measure that confound at least as much as they measure attribute
selection.

\subsection{Amazon--Walmart}

Amazon--Walmart is the largest and most attribute-rich of the three datasets (13
common attributes, 88{,}296 candidate pairs), and it is also the only dataset where
blocking itself is imperfect on the current pipeline (91.0\% pair completeness,
losing 103 of 1{,}143 true matches before the LLM is ever consulted). The one current
Step-1.5 run available reaches 0.8592 candidate-level F1 but only 0.8187 end-to-end
F1, because of that blocking loss. The learned selector keeps all 5 of the allowed
maximum attributes for every one of the 88{,}296 pairs (selection size min/mean/max
all equal 5.0). This is not a property of the selector but of the selection policy:
normalizing importance across 13 attributes spreads the mass thinly enough that the
cumulative-importance target of 0.8 cannot be met within the 5-attribute cap. Using
the LLM's own profiled importance vectors as a best case, the top 5 attributes carry
on average 0.612 of the mass and never more than 0.789, and reaching 0.8 would take
7.7 attributes on average. The cap therefore binds unconditionally, and the transfer
stage is only choosing \emph{which} five attributes, never how many. Whatever this
run measures, it is not adaptive selection in the sense Step 1.5 intends, and the
result should not be read as evidence about the method on wide schemas. No
current-pipeline Step 1 run exists for this dataset (see
Section~\ref{sec:limitations}), so there is no like-for-like global-condensation
comparison point here either.

\subsection{Cross-Dataset Synthesis}

Five observations hold across the datasets examined.

First, attribute selection is worth doing. On the one dataset with a no-selection
baseline, both mechanisms beat matching on all attributes by a wide margin
(0.7849 F1 versus 0.9231 and 0.9537), and both do so while spending fewer tokens.
Whatever the relative merits of the two mechanisms, the premise underlying both is
supported.

Second, global supervised condensation (Step 1) is competitive with or better than
the best Step-1.5 sampler on both datasets where both exist, with a far smaller
attribute footprint (2 attributes versus 4--5) and 12--18\% fewer matching tokens.
On Fodors--Zagat it leads on F1 (0.9537 versus 0.9231); on DBLP--ACM the best
Step-1.5 sampler edges ahead (0.9605 versus 0.9515) but at higher token cost. For
these two datasets a single cheaply-trained global subset captures most of what
pair-adaptive selection is trying to achieve --- but see
Section~\ref{sec:supervision}, which shows this comparison is not like-for-like.

Third, the only sampler effect that survives repeated seeds is a negative one.
Importance-based sampling is reliably the worst sampler on Fodors--Zagat (mean
0.8356 versus 0.8843--0.9222, significant against all four alternatives), and since
its profile selection is deterministic this is not one bad draw. Its apparent win on
DBLP--ACM (0.9605 versus 0.9066--0.9235) rests on single-seed runs whose reseeding
spread is wider than the gap being claimed (Section~\ref{sec:dblp-seeds}), so it
cannot currently be set against the Fodors--Zagat result. No sampler can be
recommended as a default, and the cross-dataset reversal that earlier drafts reported
is not established.

Fourth, the repeated-seed evidence changes what the sampler comparison is a
comparison \emph{of}. Among the four non-losing samplers on Fodors--Zagat the mean
differences are not statistically distinguishable, but the variances differ by an
order of magnitude, and they split precisely along whether the sampler draws its
profile pairs or derives them. This suggests the property worth designing for is
stability rather than mean accuracy, and it is a property the current
similarity-stratified default does not have. The DBLP--ACM \texttt{random} sweep is
consistent with this: SD 0.0155 over five seeds, against replicate noise of 0.0023.

Fifth, part of what the sampler comparison has been measuring is not attribute
selection. Because the selected subset is passed to the matcher in predicted-importance
order, reseeding permutes the prompt as well as changing the subset, and on DBLP--ACM
permutation alone flips 5.44\% of answers against 0.10\% for an identical prompt
(Section~\ref{sec:ordering}). Until field order is fixed to a canonical ordering, no
Step-1.5 comparison in this report cleanly isolates the variable it intends to test.
This is the single highest-value correction available, and it is a small one: sorting
the selected attributes into table-column order before condensation would remove the
confound without touching the selection logic.

\section{Limitations and Missing Data}
\label{sec:limitations}

\begin{enumerate}
\item \textbf{No current-pipeline Step 1 runs for manual, heuristic, or LLM-guided
strategies.} Every current-pipeline Step 1 run available uses the supervised
strategy; runs for the other three strategies either do not exist, or exist only
from a superseded pipeline configuration (different blocking hyperparameters,
untracked LLM model). The methodology in Section~\ref{sec:step1} for those three
strategies is accurate to the code, but no current-pipeline accuracy/token numbers
can be reported for them here. This gap is more consequential than it first appears:
\emph{manual} and \emph{heuristic} are the only label-free Step-1 strategies, and
therefore the only fair comparison points for Step 1.5 (Section~\ref{sec:supervision}).
Without them, every Step-1-versus-Step-1.5 result in this report pits a
label-consuming method against a label-free one.
\item \textbf{The Step-1 labeled sample is not held out from evaluation.} The 50
positive pairs drawn for supervised and LLM-guided selection remain in the set the
run is scored against, so those strategies are partly evaluated on pairs used to fit
them --- 45.5\% of the ground truth on Fodors--Zagat, 2.2\% on DBLP--ACM. What is
fitted is a global attribute ranking rather than a per-pair decision, so the effect is
indirect, but the Step-1 figures should be read as mildly optimistic.
\item \textbf{Amazon--Walmart has no current-pipeline Step 1 run at all}, and only
one of the five Step-1.5 samplers (similarity\_stratified) has been run on the
current embedding model and LLM. A full five-sampler sweep and a Step 1 baseline
for this dataset would substantially strengthen the cross-dataset comparison in
Section~\ref{sec:results}.
\item \textbf{Only Fodors--Zagat has a no-selection baseline}
(\texttt{--mode full}). DBLP--ACM and Amazon--Walmart have never been run without
attribute selection, so the absolute benefit of selection established in
Table~\ref{tab:fodors-baseline} rests on a single dataset.
\item \textbf{The DBLP--ACM reseeding sweep is only one-fifth complete.} Of the
thirteen runs attempted, five succeeded (\texttt{random}, seeds 42--46) and eight
failed on API exhaustion (Section~\ref{sec:dblp-seeds}). The four remaining samplers
have no repeated-seed evidence, so the ranking in Table~\ref{tab:step15-results} --
including the headline result that importance-based sampling wins on this dataset --
remains provisional, and is now known to be smaller than the reseeding spread of the
one sampler that was measured. Amazon--Walmart has no repeated-seed runs at all.
\item \textbf{Failed LLM calls are silently scored as non-matches.} The matcher
records an \texttt{api\_errors} count but \texttt{code/main.py} neither aborts nor
marks the run as invalid, so a run whose every call failed still writes a complete
\texttt{run\_summary.json} reporting precision, recall and F1. The eight aborted
DBLP--ACM runs, and a later Fodors--Zagat Step-1 run, were identified only by
inspecting \texttt{api\_errors} after the fact. The sweep driver now discards such
runs and halts, but a direct pipeline invocation still does not, so any manually
launched run must be checked before its numbers are used.
\item \textbf{Prompt field order is uncontrolled across all Step-1.5 results.}
Selected attributes reached the matcher in predicted-importance order, which varies
with the profile sample, so every Step-1.5 run varied field order alongside field
choice. On DBLP--ACM the order effect is the larger of the two
(Section~\ref{sec:ordering}). It has not been quantified on Fodors--Zagat or
Amazon--Walmart, and every Step-1.5 number in this report is affected. The pipeline
now imposes a canonical order, but no reported run was produced under it.
\item \textbf{Six runs per sampler is a small sample for a variance claim.} The
standard deviations in Table~\ref{tab:fodors-seeds} are estimated from six
observations, and three of the five samplers are deterministic in their sampling,
so their spread reflects only matcher-LLM non-determinism and is not a sampling
variance at all. The conclusion that stochastic samplers are less stable is
supported by an order-of-magnitude gap rather than by a precise estimate of either
quantity.
\item \textbf{The Step-1.5 selection policy is not scale-free and saturates on wide
schemas.} On Amazon--Walmart the cumulative-importance threshold of 0.8 is
unreachable within the 5-attribute cap: using the LLM's own profiled importance
vectors, the top 5 of 13 attributes carry on average only 0.612 of the importance
mass (maximum 0.789), so reaching 0.8 would require 7.7 attributes on average. The
cap therefore binds for 100\% of pairs regardless of what the transfer stage
predicts, which explains the constant selection size of exactly 5.0 observed in
Section~\ref{sec:results}. This is a property of normalizing importance over a
13-attribute schema, not evidence about the selector's quality, and it means the
Amazon--Walmart run does not currently test adaptive selection in any meaningful
sense.
\item \textbf{Token savings do not follow attribute savings.} On Fodors--Zagat,
Step 1.5 drops 22\% of attributes per pair but only 3.8\% of tokens, because the
selection policy ranks by importance without regard to field length and the
dropped fields are the shortest ones. Any efficiency claim for Step 1.5 needs to be
stated in tokens rather than attribute counts.
\item \textbf{Run-to-run non-determinism.} Two runs of DBLP--ACM's diversity
sampler under an identical configuration (the runs timestamped
\texttt{20260612\_153526} and \texttt{20260615\_123251}, the latter being the one
used in Table~\ref{tab:step15-results}) produced F1 of 0.9087 and 0.9162
respectively. This drift is now partly explained: the two runs predate the fixed
blocking seed, so they differ in prompt field order as well as in profile sample.
With configuration and seed genuinely held fixed the matcher is close to
deterministic (0.10--0.22\% of answers differ), so residual single-run drift should
be read as an upstream sampling artifact rather than provider-side variance.
\item \textbf{No automated leakage check exists.} The Fodors--Zagat \texttt{class}
attribute was found by manual inspection during analysis, after it had already
influenced a full generation of results. DBLP--ACM and Amazon--Walmart have since
been checked by hand and are clean, but the pipeline has no guard that would catch
a label-valued attribute in a future dataset.
\end{enumerate}

\section{Summary}

Attribute selection clearly earns its place in the pipeline: on the one dataset
where a no-selection baseline exists, both mechanisms improve F1 by 0.13--0.17
while reducing token spend. The comparison between the two mechanisms is less
favourable to the thesis's central contribution than the first generation of runs
suggested. Step 1.5's pair-adaptive selection is the more sophisticated design, but
a cheaply-trained Step 1 global condensation using only 2 attributes beats it on
Fodors--Zagat and trails it only narrowly on DBLP--ACM, in both cases at
substantially lower token cost.

That comparison, however, is not between peers. Step 1's supervised strategy is fitted
on 50 gold matches --- 45.5\% of the entire Fodors--Zagat answer key --- while Step 1.5
consumes no labels at all (Section~\ref{sec:supervision}). A method given half the
ground truth is expected to win; the informative quantity is how small the remaining
gap is without any supervision. Settling this requires the two label-free Step-1
strategies, which have never been run on the current pipeline.

The sampler comparison, which was the strongest apparent result in the earlier runs,
has shrunk under each round of scrutiny. Removing a leaking attribute from
Fodors--Zagat reversed the ranking; sweeping six runs per sampler then showed that
most of the remaining ordering was not real, since among the four non-losing samplers
the mean differences are not statistically distinguishable; and reseeding DBLP--ACM
has now shown that its ranking, too, rests on differences smaller than the run-to-run
spread. What survives is one negative result --- importance-based sampling is
reliably worst on Fodors--Zagat --- and one positive one: the samplers separate
sharply by variance rather than by mean, with deterministic samplers varying by
0.002--0.006 F1 across runs against 0.021--0.043 for the two that resample. Stability,
not mean accuracy, is the dimension on which the samplers demonstrably differ.

The reseeding work also surfaced two defects in the measurement apparatus rather than
in the method. Selected attributes were handed to the matcher in predicted-importance
order, so every Step-1.5 run varied the prompt's field order alongside its field
content; on DBLP--ACM, permuting the order while holding the attribute set fixed
changes 5.44\% of the matcher's answers, against 0.10\% for an identical prompt. And
the LLM client declared no token cap, so every request reserved 65{,}535 tokens of
budget against an actual consumption of one, aborting runs while credit remained. Both
have been fixed; neither result was available to the runs reported here.

\section{Current Status and Remaining Work}
\label{sec:status}

The findings above are stable and can be relied on: the Fodors--Zagat label leak and
its effect, the variance-versus-mean structure of the sampler comparison, the
prompt-order confound and its magnitude, the selection-policy saturation on wide
schemas, and the no-selection baseline showing that attribute selection is worth
doing. What is \emph{not} settled is the sampler ranking on either dataset.

Two code corrections have been made since the runs reported here and are covered by
regression tests, but no experiment has yet been run under them:

\begin{enumerate}
\item \textbf{Canonical prompt field order} in both selection stages
(Section~\ref{sec:ordering}). This invalidates cross-generation comparison: no
pre-fix run may be compared against a post-fix run, so the Step-1.5 tables need
regenerating in full rather than extending.
\item \textbf{Explicit token caps} on all LLM calls (Section~\ref{sec:maxtokens}),
removing the spurious budget failures that destroyed nine runs.
\end{enumerate}

Remaining experiments, in priority order, with costs measured from logged token usage
at the empirically observed rate of \$0.31 per million tokens:

\begin{center}
\small
\begin{tabular}{lrrr}
\toprule
Work item & Tokens & Est.\ cost & Runtime \\
\midrule
Step-1 \emph{manual} and \emph{heuristic}, both datasets & 9\,M & \$2.79 & 1.1\,h \\
Fodors--Zagat Step-1 baseline (re-run) & 0.7\,M & \$0.22 & 4\,min \\
Fodors--Zagat sampler sweep, 5 samplers $\times$ 5 seeds & 20\,M & \$6.19 & 1.7\,h \\
DBLP--ACM \texttt{full} and Step-1 baselines & 7.5\,M & \$2.32 & 55\,min \\
DBLP--ACM sampler sweep, 5 samplers $\times$ 5 seeds & 98\,M & \$30.35 & 8.2\,h \\
\midrule
\textbf{Total to complete the comparison} & \textbf{135\,M} & \textbf{\$41.88} & \textbf{12\,h} \\
\midrule
Amazon--Walmart, single run (blocked) & 60\,M & \$18.58 & 4.3\,h \\
\bottomrule
\end{tabular}
\end{center}

Amazon--Walmart is listed separately because it is blocked on the selection policy,
not on budget: until the cumulative-importance threshold is made scale-free, a run on
its thirteen-attribute schema saturates at \texttt{max\_k} for 100\% of pairs and
tests nothing (Section~\ref{sec:limitations}). Spending on it before that fix would
produce another uninterpretable result.

Work is currently halted for a non-technical reason: the OpenRouter account has been
exhausted. The balance stands at \$95.00 provisioned against \$95.25 consumed, leaving
it \$0.25 overdrawn, and all API calls now return HTTP 402. Completing the outstanding
comparison requires approximately \$42 of additional credit; including the
Amazon--Walmart run once its policy defect is fixed brings the total to roughly \$61.

\end{document}
